%======================================================================
%  Focused experimental design
%======================================================================

\section{Focused experimental design}
\label{sec:methodology}

\subsection{Objective and research questions}

The working objective is to determine how the two local memory systems
of a Xeon Max socket can jointly serve LLM inference.  In Flat mode,
software can place weights and KV blocks in HBM or DDR5 and can schedule
work against both tiers.  In Cache mode, HBM is a hardware-managed cache
for DDR5, so the same object-level placement policies cannot be applied.
The questions in \tabref{tab:research-questions} are provisional until
the literature review and platform configuration have been checked.
The September 2026 CRESCO8 check found Cache mode.  Flat-mode
experiments therefore require an administrator-provided configuration;
the Cache-mode question remains executable on the observed nodes.

\begin{table}[htbp]
\centering
\small
\tabsetup
\caption{Provisional research questions and testable predictions.  The
predictions are expectations, not conclusions.}
\label{tab:research-questions}
\begin{tabularx}{\textwidth}{@{}lYY@{}}
\toprule
\headrow
 & \textbf{Question} & \textbf{Prediction} \\
\midrule
RQ1
 & Which inference phases and memory objects demand HBM bandwidth, DDR
   bandwidth, or capacity on one Max 9480 socket?
 & Weight traffic dominates some low-batch decode regimes, while KV
   traffic grows with context and concurrency.  Neither tier's STREAM
   ratio alone predicts end-to-end performance. \\
RQ2
 & In Flat mode, can concurrent HBM and DDR access outperform the best
   single-tier or static split policy at the same total core count?
 & A split helps only when work in both tiers overlaps enough to offset
   scheduling and reduction overhead; the result depends on traffic,
   bandwidth saturation and thread allocation. \\
RQ3
 & How should weights and KV blocks be placed as model size, context,
   concurrency and HBM capacity change?
 & Keeping repeatedly streamed weights in HBM may beat placing the
   entire KV pool there; the best allocation can change as KV traffic
   grows.  Adaptive placement must beat tuned static baselines. \\
\addlinespace[2pt]
RQ4
 & How does hardware Cache mode compare with Flat-mode policies for
   representative working sets, if both modes are available?
 & Cache mode may help when reuse fits its effective cache capacity;
   Flat mode may do better when software can divide bandwidth and
   capacity explicitly.  No ordering is assumed. \\
\bottomrule
\end{tabularx}
\end{table}

The current vLLM CPU-tiered-KV branch is an experimental implementation
for RQ2 and RQ3, not a result that establishes either prediction.  Its
weight-in-HBM, block-level KV split, feedback allocator and parallel
attention scheduler require separate ablations.  Energy per token is a
conditional outcome until a measurement interface has been validated.
Multi-node inference and GPU comparison are outside the minimum scope.


\subsection{Evidence and the primary contrast}

Four evidence levels should remain distinct:

\begin{enumerate}
  \item a product guide establishes what a platform \emph{can} support;
  \item ENEAGRID documentation establishes what the site reported at a
        stated date;
  \item the system manifest establishes what an allocated node exposed;
        and
  \item repeated measurements establish how a specified workload
        behaved.
\end{enumerate}

The first controlled contrast holds the processor, socket, cores,
executable, model and workload constant while changing verified page
placement between local HBM and local DDR5 in Flat mode.  It calibrates
the tier effect; it is not the whole thesis question.  The next contrast
uses the same total core count and memory budget to compare sequential
and concurrent access to both tiers, followed by object-level placement
policies.  A Flat--Cache comparison, if possible, is a comparison of
memory-management modes rather than an isolated tier effect.

\subsection{Experimental sequence}

The work should advance \emph{only} after the preceding stage is
stable, in the dependency order drawn in
\figref{fig:experimental-sequence}.

\begin{figure}[htbp]
\centering
\resizebox{\textwidth}{!}{%
\begin{tikzpicture}[
  font=\small,
  stage/.style={draw=black!60, fill=neutralcol, rounded corners=2pt,
                text width=3.0cm, minimum height=1.15cm, align=center},
  primary/.style={stage, draw=accent, very thick, fill=accent!6},
  optional/.style={draw=black!45, dashed, fill=white, rounded corners=2pt,
                   text width=3.5cm, minimum height=0.95cm, align=center},
  flow/.style={-{Stealth}, thick, black!60}
]
  \node[stage] (s0) at (0,0)
    {\textbf{0. Manifest}\\topology, modes, software};
  \node[stage] (s1) at (3.7,0)
    {\textbf{1. Memory baseline}\\bandwidth, latency, locality};
  \node[primary] (s2) at (7.4,0)
    {\textbf{2. Flat calibration}\\HBM, DDR and concurrent reads};
  \node[stage] (s3) at (11.1,0)
    {\textbf{3. Flat placement}\\weights, KV and capacity};
  \node[optional] (s4) at (7.4,-2.0)
    {\textbf{Cache branch}\\working set and locality};

  \draw[flow] (s0) -- (s1);
  \draw[flow] (s1) -- (s2);
  \draw[flow] (s2) -- (s3);
  \draw[flow] (s1) -- (s4);
\end{tikzpicture}%
}
\caption{Proposed experimental dependency order.  Author's design.}
\label{fig:experimental-sequence}
\end{figure}

\begin{description}
  \item[Stage 0: system manifest.] Record the items in
  \tabref{tab:platform-record}, identify active memory and clustering
  modes, and request a Flat-mode block if available.  Only in Flat mode
  identify separate HBM and DDR NUMA nodes.

  \item[Stage 1: memory characterization.] Measure local and remote read,
  copy and triad bandwidth; unloaded and loaded latency; core-count
  scaling; and concurrent HBM--DDR traffic.  Establish saturation and
  run-to-run variability before fitting a performance model.

  \item[Stage 2: inference tier calibration, when Flat is available.]
  Freeze runtime commits,
  compiler and build.  Separate prefill from decode and compare verified
  local HBM and DDR placement.  Measure serial and concurrent access to
  both tiers with the same total core count and byte traffic.

  \item[Stage 3: object placement and capacity, when Flat is available.]
  Compare all-HBM where
  feasible, all-DDR, weights-first, KV-first and mixed-KV policies as
  context and concurrency grow.  Evaluate the existing vLLM branch
  against static allocation and attention-scheduling ablations, with
  residency and output checks.  Keep any \lcpp{} result as a separately
  named runtime baseline.

  \item[Cache-mode branch.] Measure the Cache-mode
  working-set transition on the observed nodes.  If Flat is also
  available under comparable conditions, compare modes explicitly.
  Without Flat, report Cache-mode results and leave RQ2--RQ3
  unanswered for CRESCO8 rather than implying tier control.
\end{description}

\subsection{Selecting factors without a Cartesian explosion}

The experiment matrix should be chosen from \textbf{hypotheses and
pilot data}, not by taking every combination of model, precision,
length, batch, core count and topology.  \tabref{tab:experimental-factors} lists the
factors that survive that test and states what each one is for.

\begin{table}[htbp]
\centering
\small
\tabsetup
\caption{Factors and the role of each in the experimental design.}
\label{tab:experimental-factors}
\begin{tabularx}{\textwidth}{@{}lYY@{}}
\toprule
\headrow
\textbf{Factor} & \textbf{Primary levels} & \textbf{Purpose} \\
\midrule
Memory placement
 & Local HBM, local DDR5, weights-first, KV-first and mixed KV
 & Isolate tier effects and test use of both bandwidth sources. \\
Working-set regime
 & Comfortably below, near, and above usable local HBM capacity
 & Locate a capacity boundary using measured resident memory rather than
   parameter count alone. \\
Precision
 & BF16 reference and at least one practical quantized weight format;
   KV precision recorded independently
 & Test whether reduced traffic, packing and kernel coverage change the
   HBM benefit. \\
Workload shape
 & Short and long prompt; short and sustained generation; low and
   moderate batch/concurrency
 & Separate prefill, weight-streaming decode and KV-dominant regimes. \\
Core/topology
 & Partial and full local socket; fixed total cores with varied HBM/DDR
   attention split; SNC4 only if available
 & Identify saturation, overlap and scheduling overhead. \\
Cache state
 & Cold-start/model-load and warmed steady state reported separately
 & Prevent file I/O or page-cache state from contaminating inference
   measurements. \\
\bottomrule
\end{tabularx}
\end{table}

Models should be selected by memory regime and architecture rather than
popularity alone.  The design should include a model/format that fits
comfortably in local HBM and one configuration close to the usable
capacity boundary.  At least one grouped-query model is useful because
its KV growth differs from full multi-head attention.  If the boundary
cannot be crossed with available models, context, concurrency or a
controlled allocation can be used to change the live set, provided the
mechanism does not alter the computation being compared.

Pilot runs determine the final numeric levels.  A sparse matrix of
representative points followed by focused sweeps around observed knees
is more informative than a large under-replicated grid.

\subsection{Runtime and correctness controls}

Use named, frozen commits of the experimental vLLM CPU branch for mixed
KV placement and of its unmodified base for runtime comparisons.  The
branch's present allocator assigns KV blocks to HBM or DDR for their
lifetime, and its decode path can process tier-local runs concurrently.
It does not migrate existing blocks.  Compare policies within the same
commit when possible so that allocator, kernel and model changes do not
confound the result.  A \lcpp{} run can calibrate broader CPU inference,
but its token rates are not a controlled baseline for a vLLM policy.

Performance comparisons are meaningful only when the computation is
equivalent.  Record the model revision and tokenizer, use fixed prompts
and seeds where supported, and validate token sequences or task-level
quality when precision changes.  If different kernels or quantization
formats change model quality, report that difference rather than
treating them as interchangeable.

\subsection{Metrics and explanatory measurements}

The metric definitions, equations and validity criteria are consolidated
in \secref{sec:metrics}; they are not repeated here.  The primary report
contains direct prompt-evaluation time, TTFT, the inter-token-latency
distribution, TPOT and output throughput at the declared concurrency.
Peak and steady resident memory are decomposed by object and tier where
the runtime and operating system expose that information.  Model loading
and tokenization remain outside the phase timings unless a separately
declared end-to-end experiment includes them.

The explanatory counter set is the smallest one able to test the proposed
mechanism: HBM and DDR traffic, UPI or remote traffic, achieved frequency,
cache misses and relevant matrix-engine activity.  Counter names and
availability are processor- and permission-specific.  A counter is
omitted rather than assigned a meaning that its scope does not support.
Throughput and reciprocal latency from the same synchronous batch are two
presentations of one observation, not independent confirmation.

\subsection{Run protocol and statistical analysis}

Each reported comparison should follow the same protocol:

\begin{enumerate}
  \item reserve an exclusive node and archive the system manifest;
  \item pin physical cores and memory explicitly, with SMT treated as a
        separate factor;
  \item verify page residency after model initialization and during the
        run, not only after allocation;
  \item separate model loading and warm-up from steady-state timing,
        unless they are part of the declared metric;
  \item counterbalance or randomize policy run order to reduce
        thermal and temporal bias;
  \item retain every raw sample, failure and outlier flag together with
        the command and environment;
  \item determine replication from pilot variability, using paired runs
        on the same node where possible; and
  \item repeat a subset on more than one HBM node before generalizing to
        the partition.
\end{enumerate}

Report absolute values and paired effect sizes, not only normalized bar
charts.  For primary comparisons, provide the median and interquartile
range together with a 95\% confidence interval appropriate to the
sampling design.  Tail percentiles require enough generated tokens or
requests to estimate them; p99 from a handful of observations is not a
tail-latency result.

The independent unit for a placement effect is a randomized or
counterbalanced policy run pair on a named node, not an individual
token.  Token latencies within one request share the same process,
machine state and autoregressive history and must not be counted as
independent replicates.  Confidence intervals therefore resample paired
runs (and, if the within-run distribution is analysed, preserve its
request-level grouping) rather than pooling every token into an
artificially large sample.

\subsection{Energy, if the boundary can be defended}

RAPL estimates processor-domain energy; XClarity/INA226 and a rack or
node power meter may cover different hardware.  Before reporting
joules/token, validate access, units, update rate, counter wraparound and
the components included in each source.  Integrate energy over the same
timed region as performance.  Gross joules/request and
joules/output-token should be reported; if idle energy is subtracted,
also report the gross value and the subtraction procedure.  Facility
cooling claims require a measured facility boundary and are outside a
node-level study.

\subsection{Threats to validity}

\tabref{tab:validity} pairs each class of threat with the design
feature that answers it.  The pairing matters more than the list: a
threat with no corresponding response is one the study cannot claim to
have addressed.

\begin{table}[htbp]
\centering
\small
\tabsetup
\caption{Main validity threats and design responses.}
\label{tab:validity}
\begin{tabularx}{\textwidth}{@{}lYY@{}}
\toprule
\headrow
\textbf{Validity} & \textbf{Threat} & \textbf{Response} \\
\midrule
Internal
 & Frequency drift, first touch, page migration, remote access, warm
   cache state or shared-node interference is mistaken for an HBM effect.
 & Exclusive nodes, paired order, explicit binding, residency checks and
   frequency/traffic measurements. \\
Construct
 & STREAM is treated as LLM bandwidth; file size as live memory; average
   TPOT as tail latency; RAPL as whole-node energy.
 & Define each metric and boundary, measure the live set, and use
   counters only to test a stated mechanism. \\
External
 & One model, precision, node or runtime is generalized to CPU inference
   as a whole.
 & Cover distinct memory regimes, repeat on multiple nodes, and state the
   platform/runtime boundary of every conclusion. \\
Statistical conclusion
 & Few repetitions, selected best runs or unreported failures inflate an
   effect.
 & Retain raw samples, predeclare primary comparisons and report
   uncertainty with absolute values. \\
Reproducibility
 & Modules, firmware, models and scheduler inventory change over time.
 & Archive immutable source/model revisions, build metadata, commands
   and dated system manifests. \\
\bottomrule
\end{tabularx}
\end{table}


\subsection{When an implementation contribution is justified}

The experimental vLLM branch already contains an object-level placement
implementation.  The present task is to determine whether its mixed KV
allocation and concurrent attention improve measured performance over
strong static baselines, and under which workloads.  Further policy
changes should target a measured inefficiency rather than be assumed
beneficial.

Compare all-local placement when it fits, weights-first and KV-first
tiering, static mixed fractions, the adaptive allocator, and the tier
attention scheduler both enabled and disabled.  Include placement and
reduction overhead, remote traffic, memory use and output quality.
Claim novelty only after the revised literature review and controlled
measurements establish what differs from prior tiered-attention work.
